\subsection{Preparación de los datos en los experimentos iniciales}

Los archivos \texttt{ridge.py}, \texttt{OLS\_10.py} y \texttt{ridge\_10.py} leen \texttt{datos.csv}. La variable de respuesta es la generación eléctrica, identificada como \texttt{Generacion (MWh)}. A partir de \texttt{Timestamp} se extraen componentes del calendario y se construyen cinco predictores: hora del día con fracción de minutos, día del año, año, mes y día de la semana. El área seleccionada corresponde al primer valor distinto que aparece en la columna \texttt{Area}.

El filtro comprende desde el 1 de enero de 2022 hasta el 31 de diciembre de 2022 a las 00:00, inclusive; por tanto, no equivale necesariamente a incluir todas las horas del último día del año. La columna auxiliar de fecha se reconstruye con año, mes, día y hora, sin incorporar los minutos. Después se conserva el primer 70\,\% de las filas disponibles para generar las ventanas. El 30\,\% restante no se evalúa en estos tres programas. Las versiones de diez pasos ordenan explícitamente por la fecha auxiliar; la versión Ridge de un paso conserva el orden de lectura del archivo. En consecuencia, interpretar los pasos como horas consecutivas requiere verificar previamente la frecuencia y el orden de los registros.
